\section{Translation: the label model as a coarsening problem}
\label{sec:translation}

\subsection{The data-generating process}

Consider a binary classification task with true label $Y \in \{0,1\}$
and class prior $\pi = \Prob\{Y = 1\}$; the multiclass case is a
direct extension. There are $m$ labeling functions. For example $i$,
LF $j$ emits a vote $\lambda_{ij} \in \{0, 1, \varnothing\}$, where
$\varnothing$ denotes \emph{abstention}. Each LF $j$ has a
\emph{coverage} $\beta_j = \Prob\{\lambda_{ij} \neq \varnothing\}$,
the probability it does not abstain, and an \emph{accuracy}
$\alpha_j = \Prob\{\lambda_{ij} = Y \mid \lambda_{ij} \neq
\varnothing\}$, the probability it votes correctly given that it
votes. When LF $j$ votes, it emits $Y$ with probability $\alpha_j$
and the complementary label $1 - Y$ with probability $1 - \alpha_j$.

The label model is fit to the matrix of votes $\{\lambda_{ij}\}$. Its
parameters are the per-LF accuracies $\bm{\alpha} = (\alpha_1, \ldots,
\alpha_m)$, the coverages $\bm{\beta}$, and the class prior $\pi$.
Coverages are directly observed (each is an empirical abstention
rate), so the inferential content is the recovery of $\bm{\alpha}$
and $\pi$, after which the posterior $\Prob\{Y = y \mid
\bm{\lambda}_i\}$ supplies the probabilistic training label for
example $i$.

\subsection{LF votes as a candidate set}

The votes of the $m$ LFs on example $i$ induce a \emph{candidate set}
$c_i \subseteq \{0, 1\}$: the set of labels still consistent with the
observed votes under the modeled noise. An LF that abstains
constrains nothing. An LF that votes is consistent with both labels
(it could be voting correctly or incorrectly) until its accuracy is
known, but its vote shifts the posterior mass. In the deterministic
limit, where an LF that votes is treated as ground truth, a single
confident vote collapses the candidate set to that label; a
gold-labeled example fixes the candidate set to the singleton
$\{Y_i\}$ by construction. The probabilistic label model is the
soft-coarsening relaxation of this combinatorial picture: instead of
a hard candidate set, each example carries a distribution over
$\{0,1\}$ determined by the LF votes and the accuracy parameters.

This candidate-set reading of weak supervision is itself the subject
of an established line of work, which we connect to rather than
introduce. \citet{yu2022nplm} make LFs emit candidate-label subsets
explicitly and aggregate them with a gold-free generative model; the
partial-label and superset-label learning literatures study learning
a classifier when each example carries a candidate set known to
contain the truth, with the learnability boundary controlled by an
\emph{ambiguity degree} \citep{cour2011partial}, sample-complexity
conditions for empirical-risk minimization on supersets
\citep{liu2014superset}, consistent disambiguation losses with proved
rates \citep{cabannes2020infimum}, and a mixing-matrix admissibility
condition for proper losses on partial labels
\citep{cidsueiro2012proper}. Those works fix the candidate set as
given supervision and study how to learn from it; we instead treat the
candidate set as the \emph{output} of a coarsening mechanism whose
own parameters (LF accuracies, prior) must be identified, and read
the C1--C2--C3 conditions as the admissibility statement on that
mechanism.

\subsection{The translation}

\Cref{tab:translation} states the explicit correspondence to
masked-cause inference.

\begin{table}[ht]
\centering
\small
\begin{tabular}{@{}p{0.40\linewidth} p{0.56\linewidth}@{}}
\toprule
\textbf{Masked-data framework} & \textbf{Programmatic weak supervision} \\
\midrule
Component / latent cause $K_i$ & True label $Y$ \\
Candidate set $c_i$ & Labels consistent with the LF votes for an example \\
Masking probability & The per-LF accuracy / confusion model \\
Singleton candidate set $|c_i|=1$ & A gold-labeled (hand-annotated) example \\
Non-informative (full) candidate set & An LF that abstains \\
Narrowed candidate set & A confident LF vote \\
Augmented candidate-set matrix $\tilde C$ & The matrix of LF votes
  augmented with any gold labels; full column rank gives
  identifiability \\
C1 (truth in candidate set) & True label consistent with the votes;
  holds unless LFs are confidently wrong beyond modeled noise \\
\textbf{C2} (coarsening at random): conditional on the candidate set, the masking distribution does not depend on which element is the true cause &
  Concretely: LF error depends on $Y$ only through the modeled
  confusion, not on the label-model parameters \\
C3 (parameter independence) & LF accuracy parameters distinct from
  the downstream model parameters \\
\bottomrule
\end{tabular}
\caption{Translation between the masked-cause framework
\citep{towell2026masked} and programmatic weak supervision.}
\label{tab:translation}
\end{table}

The three conditions deserve a reading in the weak-supervision
language.

\paragraph{C1, support.} The true label must remain consistent with
the LF votes. This fails exactly when an LF is confidently wrong in a
way the noise model does not allow, for example an LF whose error is
systematic on a subpopulation (a label-conditional bias the
symmetric accuracy parameter cannot represent). For ordinary noisy
LFs, C1 holds: the modeled noise already admits any vote pattern.

\paragraph{C2, symmetry.} LF error must depend on $Y$ only through
the modeled confusion. An LF whose error rate is itself a function of
the label-model parameters, or whose mistakes are correlated with the
quantity being estimated, breaks C2. The standard data-programming
accuracy model satisfies C2 by construction: $\Prob\{\lambda_{ij}
\mid Y\}$ is a fixed confusion table.

\paragraph{C3, parameter independence.} The LF accuracy parameters
$\bm{\alpha}$ must be distinct from the downstream classifier's
parameters. This holds in the standard pipeline, where the label
model is fit first and the downstream model second; it would fail in
a joint end-to-end scheme that ties LF reliability to downstream
performance.

\subsection{Existing methods in this language}

\paragraph{Dawid--Skene \citep{dawid1979maximum}} fit per-rater
confusion matrices and a class prior by EM. In our language: the
candidate-set coarsening is the confusion table, and EM is the
maximum-likelihood solver for the coarsened likelihood
\eqref{eq:bg-c1c2c3}. Dawid--Skene is the clear ancestor of the
label model.

\paragraph{Data programming \citep{ratner2016data}} fits a
factor-graph label model and, in the independent case, recovers LF
accuracies from agreement statistics. This is identifiability under
C1--C2--C3 with conditionally independent LFs (\cref{sec:identifiability}).

\paragraph{Snorkel \citep{ratner2017snorkel}} is the system built on
data programming; the multi-task extension \citep{ratner2019training}
handles structured LF dependence with a known dependency graph. In
our language, a known dependency graph is a structural assumption
that repairs the rank deficit C2 violations would otherwise cause.

\paragraph{Triplet method \citep{fu2020fast}} recovers LF accuracies
in closed form from third-order agreement rates of conditionally
independent LFs. This is the constructive proof of T2 below.

\paragraph{Crowdsourcing \citep{karger2011iterative}} solves the same
estimation problem with workers in place of LFs and message-passing
in place of EM. The masked-cause frame covers both: a worker and an
LF are both noisy coarsenings of the true label.

These methods are not in conflict; they are different solvers and
parametrizations of one coarsening problem. The framework states
precisely when each is identifiable and what auxiliary information
(gold labels, a dependency graph, a better-than-random assumption)
each implicitly requires.
